\section{Discusión}\label{sec:discusion}

\subsection{Comparación entre estrategias de preprocesamiento}

Desde un punto de vista estrictamente numérico, la estrategia
técnico-estadística (V2) obtuvo métricas superiores. Mientras que en la
estrategia clínico-estadística (V1) los mejores modelos, Random Forest y GBM,
alcanzaron valores de AUC próximos a 0.947--0.949 en validación y 0.939--0.942
en prueba, en la V2 estos mismos modelos alcanzaron valores cercanos a 0.978 en
validación y 0.974--0.975 en prueba.

Sin embargo, en medicina predictiva, el mayor rendimiento métrico no implica
necesariamente una mayor validez clínica. La diferencia observada en AUC sugiere
que la estrategia V2 podría beneficiarse de variables que la estrategia V1
excluyó por riesgo de \textit{Data Leakage} o por menor coherencia biológica. En
este sentido, la V2 ofrece mayor rendimiento estadístico, mientras que la V1
representa una aproximación más prudente desde el punto de vista clínico y
metodológico.

\subsection{Contraste de la hipótesis}

La hipótesis planteada proponía que el preprocesamiento con criterio clínico
mejoraría la calidad del modelo. Los resultados muestran una confirmación
parcial:

\begin{itemize}
  \item En términos de estabilidad, la estrategia V1 mantiene resultados muy
    consistentes entre la versión con CV y sin CV para los modelos de ensamble.
  \item En términos de rendimiento bruto, la V2 supera a la V1.
  \item En términos de evaluación final, el conjunto de prueba confirma la misma
    tendencia observada durante la validación, sin alterar la selección del mejor
    modelo.
\end{itemize}

Esta diferencia refuerza la necesidad de interpretar el rendimiento predictivo
junto con la procedencia y disponibilidad real de las variables. La V2
probablemente retuvo información que facilita la predicción estadística, pero
que en un entorno hospitalario real al momento del ingreso podría no estar
disponible o presentar riesgo de contaminación temporal. La V1, aunque no
maximiza el AUC, mantiene un rendimiento elevado (AUC > 0.93 en prueba) bajo un
criterio de mayor fidelidad clínica.

\subsection{Impacto del conocimiento clínico}

El impacto del conocimiento clínico se observa en la curación del conjunto de
datos:

\begin{enumerate}
  \item Imputación estratificada: Aunque la V1 utiliza un conjunto más restringido
    de variables predictoras, la combinación de imputaciones contextualizadas y
    eliminaciones estratégicas permitió que modelos como Random Forest
    y GBM alcanzaran AUC superiores a 0.94 en validación. Esto sugiere que los
    predictores conservados mantienen una alta capacidad informativa.
  \item Sensibilidad y especificidad: En V1, la especificidad se mantiene en
    valores elevados (aprox. 0.93--0.95). Este comportamiento indica una buena
    capacidad para identificar correctamente a los pacientes supervivientes, un
    aspecto relevante para reducir falsas alarmas en contextos de asignación de
    recursos críticos.
  \item Rendimiento de KNN: Las Tablas~\ref{tab:v1-cv-base},
    \ref{tab:v1-nocv-base}, \ref{tab:v2-cv-base} y
    \ref{tab:v2-nocv-base} muestran que KNN fue superado por los métodos
    basados en árboles, aun cuando la normalización fue integrada dentro del
    flujo de tidymodels. Esto valida que la relación entre las variables
    clínicas y la mortalidad no es necesariamente lineal ni depende solo de
    distancias simples.
\end{enumerate}

\subsection{Comportamiento de los modelos}

\begin{itemize}
  \item Modelos de ensamble: Random Forest (RF) y Stochastic Gradient Boosting
    (GBM) mostraron el mejor rendimiento global. GBM lideró los escenarios V1,
    mientras que Random Forest y GBM presentaron rendimientos prácticamente
    equivalentes en V2. Este patrón refleja su capacidad para capturar
    interacciones no lineales entre variables clínicas.
  \item Regresión logística: Este modelo alcanzó AUC superiores a 0.93 en Bio Data,
    lo que sugiere que varias de las variables clínicas seleccionadas mantienen
    una relación discriminativa estable con la mortalidad. Además, su mayor
    interpretabilidad la convierte en una referencia útil frente a modelos más
    complejos.
  \item Support Vector Machine: La incorporación de SVM permitió evaluar un modelo
    de margen máximo plenamente integrado en el flujo tidymodels. Su rendimiento
    fue competitivo, especialmente en V2, aunque no superó de forma consistente
    a los modelos de ensamble.
  \item Decision Tree (rpart): A diferencia de experimentos previos, el árbol de
    decisión simple mantuvo un rendimiento discriminativo razonable en los
    escenarios actualizados. Sin embargo, continuó por debajo de Random Forest y
    GBM, lo que subraya la ventaja de los métodos de ensamble en problemas
    biomédicos con relaciones no lineales y variables heterogéneas.
\end{itemize}
